\begin{table}[t]
\caption{Разведочные контрасты качества прямого решателя. Разности и интервалы указаны в процентных пунктах; положительные значения благоприятствуют D.}
\label{tab:direct-contrasts}
\centering\small
\begin{tabular}{@{}lrrr@{}}
\toprule
Контраст & Задачи & Разность & 95\% интервал\\
\midrule
D -- SN & 980 & $+0.27$ & $[-0.88, 1.43]$\\
D -- SR & 979 & $+0.95$ & $[-0.17, 2.11]$\\
D -- MN & 966 & $+4.59$ & $[2.97, 6.18]$\\
D -- MR & 970 & $+3.20$ & $[1.82, 4.64]$\\
\bottomrule
\end{tabular}
\par\smallskip\noindent Интервалы получены повторной выборкой задач. Эти апостериорные контрасты не входят в семейство из четырех тестов с поправкой Холма и предназначены для оценивания, а не для образования нового подтверждающего семейства.
\end{table}
